\chapter{Communication: Serialization and Compression}
\label{ch:communication}

\section{Message Serialization}
\label{sec:serialization}

In a federated learning pipeline, the dominant network payload consists of model weight
tensors exchanged between edge nodes and the central coordinator during each training
round. The choice of serialization format directly impacts encoding and decoding
latency on resource-constrained edge hardware, wire-format payload size, and
cross-platform interoperability across heterogeneous node environments
\cite{mcmahan2017communication}.

Three serialization formats are considered:

\begin{itemize}
   \item \textbf{JSON}: A human-readable, text-based format suitable for control
   messages and configuration payloads \cite{rfc8259}. It is generally inefficient
   for dense numerical arrays because values are encoded as text rather than fixed-width
   binary primitives \cite{ahmad2022arrowflight, gudla2024cbl, blyth2018proio, 
   viotti2022benchmark}.

    \item \textbf{Protocol Buffers (Protobuf)}: A language-neutral binary serialization
    framework developed by Google \cite{protobuf_overview, grpc2016whitepaper}. It uses
    compact binary encoding, with variable-length representations for certain primitive
    types and fixed-width encoding for floating-point values.

   \item \textbf{FlatBuffers}: A binary format designed for direct access to
   serialized data without full deserialization, making it suitable for
   latency-sensitive or CPU-constrained settings \cite{flatbuffers_docs,
   flatbuffers_whitepaper}. Existing benchmarks show that it does not inherently
   guarantee smaller payloads than other binary serialization formats
   \cite{viotti2022benchmark, protobuf_overview}.
\end{itemize}

\begin{table}[H]
    \centering
    \begin{tabular}{|p{0.22\textwidth}|p{0.22\textwidth}|p{0.22\textwidth}|p{0.22\textwidth}|}
        \hline
        \textbf{Criterion} & \textbf{JSON} & \textbf{Protobuf} & \textbf{FlatBuffers} \\
        \hline
        Payload Size & Large (text) & Compact (binary) & Compact (binary) \\
        \hline
        Serialization Speed & Slow & Fast & Fast \\
        \hline
        Deserialization & Slow (parsing) & Fast & Zero-copy access \\
        \hline
        Schema Enforcement & None & Strong (.proto) & Strong (.fbs) \\
        \hline
        Human Readability & High & None & None \\
        \hline
    \end{tabular}
    \caption{Comparison of serialization formats for federated model payloads}
    \label{tab:serial_comparison}
\end{table}

Overall, binary serialization formats such as Protobuf and FlatBuffers are generally
preferred for high-volume machine-to-machine communication because they reduce payload
size and parsing overhead relative to text-based representations. JSON remains valuable
primarily for readability, debugging, and human-facing control messages
\cite{rfc8259, viotti2022benchmark, protobuf_overview, flatbuffers_whitepaper}.

\section{Message Transport}
\label{sec:transport}

Model update exchange in distributed and federated learning is generally treated as
loss-sensitive, since missing or corrupted updates may reduce aggregation fidelity,
introduce staleness, or interfere with synchronization across training rounds.
Accordingly, most practical distributed learning systems assume reliable communication,
and transport mechanisms with ordered, loss-recovering delivery are commonly preferred
for model exchange workloads \cite{yu2019distributed, mcmahan2017communication}.

Transmission Control Protocol (TCP) provides reliable, ordered delivery with 
retransmission, congestion control, and flow control \cite{rfc793}. By contrast, 
User Datagram Protocol (UDP) offers lower protocol overhead but does not guarantee 
delivery, ordering, or retransmission \cite{rfc768}. Accordingly, TCP is commonly 
favored for model and gradient exchange, whereas UDP is more suitable for auxiliary 
signaling or lightweight telemetry where occasional packet loss is acceptable.

A common higher-level transport framework in distributed systems is gRPC
\cite{grpc2016whitepaper}, which typically operates over HTTP/2. Its main
communication advantages include:

\begin{itemize}
    \item \textbf{HTTP/2 Multiplexing}: Multiple concurrent RPC streams can share a
    single TCP connection, reducing connection-management overhead and improving
    network utilization \cite{rfc7540, rfc9113}.

    \item \textbf{Bidirectional Streaming}: Both endpoints may exchange data over the
    same long-lived stream, which is advantageous for iterative coordination and
    repeated model exchange \cite{grpc2016whitepaper}.

    \item \textbf{Native Protobuf Integration}: Service definitions in \texttt{.proto}
    files enable strongly typed cross-language communication \cite{protobuf_overview, 
    grpc2016whitepaper}.

    \item \textbf{Compression Support}: gRPC supports message compression at the RPC
    layer, allowing transport-level payload reduction without modifying the application
    protocol \cite{grpc2016whitepaper}.
\end{itemize}

As a result, TCP-based RPC frameworks such as gRPC are widely used in distributed
machine learning and microservice-oriented systems due to their balance of reliability,
efficiency, and interoperability.

\section{Message Compression}
\label{sec:compression}

Even after binary serialization, the bulk of the payload often remains composed of
dense floating-point arrays. Communication reduction is therefore commonly addressed 
through both \emph{algorithmic compression} and \emph{byte-level compression}.

\subsection{Top-K Gradient Sparsification}
\label{subsec:compression_topk}

Top-K gradient sparsification \cite{aji2017sparse} reduces communication volume at the
algorithmic level. Rather than transmitting the full gradient vector
$\nabla W \in \mathbb{R}^P$, only the $K$ components with the largest absolute
magnitudes are retained:

\begin{equation}
    \text{Top-K}(\nabla W) = \{(i, \nabla W_i) : i \in \text{argtop-K}(|\nabla W|)\}
    \label{eq:topk}
\end{equation}

The remaining $(P - K)$ components are commonly accumulated locally in a residual
buffer and reintroduced in subsequent rounds, ensuring that omitted information 
is not permanently discarded \cite{lin2018deep_gradient, karimireddy2019errorfeedback, 
stich2018sparsified}. This approach can substantially reduce communication volume 
while preserving model convergence when appropriately configured.

\subsection{Byte-Level Compression}
\label{subsec:compression_byte}

After serialization, the payload may be further reduced using general-purpose
lossless compression algorithms. Two widely used choices are \emph{LZ4} and
\emph{Zstandard (Zstd)}.

\begin{itemize}
    \item \textbf{LZ4} \cite{lz4_reference, collet2011lz4}: A high-speed lossless
    compression algorithm optimized primarily for throughput rather than maximum
    compression ratio. Published benchmarks and official documentation report
    compression throughput in the hundreds of MB/s and decompression throughput
    reaching multiple GB/s on modern CPUs, with relatively low CPU overhead.
    Its compression ratio is typically moderate compared with more aggressive codecs.

    \item \textbf{Zstandard (Zstd)} \cite{collet2018zstandard}: A modern lossless
    compression algorithm designed to balance high throughput with stronger
    compression ratios. Across a wide range of workloads, Zstd generally achieves
    better compression than LZ4 while maintaining high decompression speed, though
    at higher compression-time CPU cost depending on the selected compression level.
\end{itemize}

\begin{table}[H]
    \centering
    \begin{tabular}{|p{0.28\textwidth}|p{0.30\textwidth}|p{0.30\textwidth}|}
        \hline
        \textbf{Metric} & \textbf{LZ4} & \textbf{Zstd} \\
        \hline
        Compression Ratio & Moderate & Higher \\
        \hline
        Compression Speed & Very High & High \\
        \hline
        Decompression Speed & Very High & Very High \\
        \hline
        CPU Usage & Low & Moderate (level-dependent) \\
        \hline
        Best Use Case & CPU-constrained settings & Bandwidth-constrained settings \\
        \hline
    \end{tabular}
    \caption{Conceptual comparison of LZ4 and Zstandard based on reported literature}
    \label{tab:compression_comparison}
\end{table}

In general, the trade-off between LZ4 and Zstd is straightforward: LZ4 prioritizes
minimal computational overhead and very high speed, whereas Zstd prioritizes stronger
payload reduction while preserving practical throughput
\cite{lz4_reference, collet2018zstandard}.

\subsection{Adaptive Compression Selection}
\label{subsec:compression_adaptive}

In heterogeneous distributed environments, the most appropriate compression algorithm
depends on system constraints such as CPU availability, memory pressure, and network
bandwidth. For this reason, prior work and practical distributed systems often motivate
\emph{adaptive compression strategies}, in which the compression method or level is
selected dynamically according to current operating conditions
\cite{collet2018zstandard, konevcny2016federated, kairouz2021advances}.

Conceptually, lightweight algorithms such as LZ4 are more suitable when minimizing
local CPU overhead is the primary objective, whereas stronger codecs such as Zstd are
more advantageous when network bandwidth is the dominant bottleneck. This reflects a
general systems trade-off between computation and communication cost.
